\section{AHSME 1968}
        \begin{enumerate}
        	\item (\textit{AHSME 1968}) Let $P$ units be the increase in circumference of a circle resulting from an increase in $\pi$ units in the diameter. Then $P$ equals:


 $\text{(A) } \frac{1}{\pi}\quad\text{(B) } \pi\quad\text{(C) } \frac{\pi^2}{2}\quad\text{(D) } \pi^2\quad\text{(E) } 2\pi$

 \textbf{Answer: }D



\textbf{Solution}

Let $d$ be the diameter of the original circle. If $d$ is increased by $\pi$, then the new circumference is $\pi d + \pi^2$. The difference in circumference is therefore $\pi d + \pi^2 - \pi d = \pi^2$ 

 Therefore, the answer is $\fbox{D}$

 Solution by VivekA

 


	\item (\textit{AHSME 1968}) The real value of $x$ such that $64^{x-1}$ divided by $4^{x-1}$ equals $256^{2x}$ is:


 $\text{(A) } -\frac{2}{3}\quad\text{(B) } -\frac{1}{3}\quad\text{(C) } 0\quad\text{(D) } \frac{1}{4}\quad\text{(E) } \frac{3}{8}$

 \textbf{Answer: }B



\textbf{Solution}

Because $64=4^3$ and $256=4^4$, we can change all of the numbers in the equation to exponents with base $4$ and solve the equation:\begin{align*}\frac{64^{x-1}}{4^{x-1}}=256^{2x} \\\frac{(4^3)^{x-1}}{4^{x-1}}=(4^4)^{2x} \\\frac{4^{3x-3}}{4^{x-1}}=4^{8x} \\4^{2x-2}=4^{8x} \\2x-2=8x \\x-1=4x \\ 3x=-1 \\x=\frac{-1}{3} \\\end{align*}


 Thus, our desired answer is $\fbox{(B) -1/3}$.

 


	\item (\textit{AHSME 1968}) A straight line passing through the point $(0,4)$ is perpendicular to the line $x-3y-7=0$. Its equation is:


 $\text{(A) } y+3x-4=0\quad \text{(B) } y+3x+4=0\quad \text{(C) } y-3x-4=0\quad \\ \text{(D) } 3y+x-12=0\quad \text{(E) } 3y-x-12=0$

 \textbf{Answer: }A



\textbf{Solution}

The original line can be adjusted to \href{https://artofproblemsolving.com/wiki/index.php?title=Standard\_form}{standard form}, in which it is $x-3y=7$. Because this line is in standard form, its slope is $\frac{-1}{-3}=\frac{1}{3}$. Thus, the slope of the perpendicular line is the negative reciprocal of this number, $-3$. Because we know that this new line passes through the point $(0,4)$, we can describe this line using \href{https://artofproblemsolving.com/wiki/index.php?title=Point-slope\_form}{point-slope form}, in which it is $y-4=-3(x-0)$, or $y+3x-4$, which is answer choice $\fbox{A}$.

 


	\item (\textit{AHSME 1968}) Define an operation $\star$ for positive real numbers as $a\star b=\frac{ab}{a+b}$. Then $4 \star (4 \star 4)$ equals:


 $\text{(A) } \frac{3}{4}\quad \text{(B) } 1\quad \text{(C) } \frac{4}{3}\quad \text{(D) } 2\quad \text{(E )} \frac{16}{3}$

 \textbf{Answer: }C



\textbf{Solution}

$4 \star 4 \star 4 = 4 \star 2 = \frac{4}{3}$

 Therefore, our answer is $\fbox{C}$

 Solution by VivekA

 


	\item (\textit{AHSME 1968}) If $f(n)=\frac{1}{3} n(n+1)(n+2)$, then $f(r)-f(r-1)$ equals:


 $\text{(A) } r(r+1)\quad \text{(B) } (r+1)(r+2)\quad \text{(C) } \frac{1}{3} r(r+1)\quad  \\ \text{(D) } \frac{1}{3} (r+1)(r+2)\quad \text{(E )} \frac{1}{3} r(r+1)(2r+1)$

 \textbf{Answer: }A



\textbf{Solution}

Plugging in the expressions for $f(r)$ and $f(r-1)$, we see that:\begin{align*}f(r)-f(r-1)&= \frac{1}{3}r(r+1)(r+2)-\frac{1}{3}(r-1)(r-1+1)(r-1+2) \\&=\frac{1}{3}[r(r+1)(r+2)-r(r-1)(r+1)] \\&=\frac{1}{3}[r(r^2+3r+2)-r(r^2-1)] \\&=\frac{1}{3}[r^3+3r^2+2r-r^3+r] \\&=\frac{1}{3}[3r^2+3r] \\&=r^2+r \\&=r(r+1), \\\end{align*}
which is answer choice $\fbox{A}$.

 


	\item (\textit{AHSME 1968}) Let side $AD$ of convex quadrilateral $ABCD$ be extended through $D$, and let side $BC$ be extended through $C$, to meet in point $E.$ Let $S$ be the degree-sum of angles $CDE$ and $DCE$, and let $S'$ represent the degree-sum of angles $BAD$ and $ABC.$ If $r=S/S'$, then:


 $\text{(A) } r=1 \text{ sometimes, } r>1 \text{ sometimes}\quad\\ \text{(B) }  r=1 \text{ sometimes, } r<1 \text{ sometimes}\quad\\ \text{(C) } 0<r<1\quad \text{(D) } r>1\quad \text{(E) } r=1$

 \textbf{Answer: }E



\textbf{Solution}


                    \begin{figure}[H]
                        \centering
                        \includegraphics[width=0.5\textwidth]{Solution/AHSME_I_1968_6_Solution_Figure_1.png}
                    \end{figure}
                    

 Because $ABCD$ is a convex \href{https://artofproblemsolving.com/wiki/index.php?title=Quadrilateral}{quadrilateral}, the sum of its interior angles is $360^{\circ}$. Thus, $S'=\measuredangle BAD + \measuredangle ABC=360^{\circ}-(\measuredangle BCD+\measuredangle ADC)$. Furthermore, because $\angle BCD$ and $\angle ADC$ are \href{https://artofproblemsolving.com/wiki/index.php?title=Supplementary}{supplementary} to $\angle DCE$ and $\angle CDE$, respectively, the four angles sum to $2*180^{\circ}=360^{\circ}$, so $\measuredangle BCD+\measuredangle ADC=360^{\circ}-(\measuredangle DCE+\measuredangle CDE)=360^{\circ}-S$. Plussing this expression for $\measuredangle BCD+\measuredangle ADC$ into the first equation, we see that $S'=360^{\circ}-(360^{\circ}-S)=S$, so $r=\frac{S}{S'}=1$, which is answer choice $\fbox{E}$.

 


	\item (\textit{AHSME 1968}) Let $O$ be the intersection point of medians $AP$ and $CQ$ of triangle $ABC.$ if $OQ$ is 3 inches, then $OP$, in inches, is:


 $\text{(A) } 3\quad \text{(B) } \frac{9}{2}\quad \text{(C) } 6\quad \text{(D) } 9\quad \text{(E) } \text{undetermined}$

 \textbf{Answer: }E



\textbf{Solution}

The fact that $OQ=3$ only tells us that $CQ=9$. There are infinitely many triangles with a median which has length 9, so we can make no statement about the length of the median $\overline{AP}$ or the segment $\overline{OP}$. Thus, $OP$ is $\fbox{(E) undetermined}$.

 


	\item (\textit{AHSME 1968}) A positive number is mistakenly divided by $6$ instead of being multiplied by $6.$ Based on the correct answer, the error thus committed, to the nearest percent, is :


 $\text{(A) } 100\quad \text{(B) } 97\quad \text{(C) } 83\quad \text{(D) } 17\quad \text{(E) } 3$

 \textbf{Answer: }B



\textbf{Solution}

Let $n$ be the original number in question. The number divided by 6 is $\frac{n}{6}$, and the number multiplied by 6 is $6n$. Because $\frac{n}{6}=\frac{1}{36}*6n$, and $\frac{1}{36}=0.02 \overline{7} \approx 0.03=3$\%, the error (which is our deviation from the correct value) is about $97$\%. Thus, our answer is $\fbox{B}$.

 


	\item (\textit{AHSME 1968}) The sum of the real values of $x$ satisfying the equality $|x+2|=2|x-2|$ is:


 $\text{(A) } \frac{1}{3}\quad \text{(B) } \frac{2}{3}\quad \text{(C) } 6\quad \text{(D) } 6\frac{1}{3}\quad \text{(E) } 6\frac{2}{3}$

 \textbf{Answer: }E



\textbf{Solution}

There are three cases: $\newline$(i) $(x+2), (x-2) \geq 0$ $\newline$(ii) $x+2 \geq 0, x-2<0$ $\newline$(iii) $(x+2), (x-2)<0$ $\newline$In case (i), we solve the equation $x+2=2(x-2)$ to get $x=6$. In case (ii), we solve the equation $x+2=2(2-x)$ to get $x=\frac{2}{3}$. Case (iii) simply produces an equation which is that of case (i) multiplied on both sides by $-1$, so it yields the same solutions and thus may be discarded. Thus, if we add the real values of $x$ which satisfy the original equation, we get $6+\frac{2}{3}=6\frac{2}{3}$, or answer choice $\fbox{E}$.

 


	\item (\textit{AHSME 1968}) Assume that, for a certain school, it is true that


 I: Some students are not honest.
II: All fraternity members are honest.


 A necessary conclusion is:


 $\text{(A) Some students are fraternity members.} \quad\\ \text{(B) Some fraternity member are not students.} \quad\\ \text{(C) Some students are not fraternity members.} \quad\\ \text{(D) No fraternity member is a student.} \quad\\ \text{(E) No student is a fraternity member.}$

 \textbf{Answer: }C



\textbf{Solution}

If some students are dishonest, we know that they must not be fraternity members, because if they were members, then they would be honest. This conclusion aligns with answer choice $\fbox{C}$.

 Choice (A) is incorrect, because if there were no fraternity members, proposition II would be \href{https://en.wikipedia.org/wiki/Vacuous\_truth}{vacuously true}, and dishonest students are still allowed to exist within the rules of the problem, satisfying propostion I.

 Choice (B) is incorrect, because dishonest students can still exist (proposition I) while some other students (the fraternity members, maybe some others) are honest (proposition II).

 Choice (D) is incorrect by the same reasoning used against choice (B).

 Choice (E) is incorrect, because it is the \href{https://artofproblemsolving.com/wiki/index.php?title=Contrapositive}{contrapositive} of choice (D), so the two are logically equivalent. (Choice (D) could be restated as "if someone is a fraternity member, that someone is not a student," and choice (E) could be restated as "if someone is a student, that someone is not a fraternity member.")

 


	\item (\textit{AHSME 1968}) If an arc of $60^{\circ}$ on circle $I$ has the same length as an arc of $45^{\circ}$ on circle $II$, the ratio of the area of circle $I$ to that of circle $II$ is:


 $\text{(A) } 16:9\quad \text{(B) } 9:16\quad \text{(C) } 4:3\quad \text{(D) } 3:4\quad \text{(E) } \text{none of these}$

 \textbf{Answer: }B



\textbf{Solution}

Let the length of the two arcs be denoted $x$. First, we acknowledge that the length of an \href{https://artofproblemsolving.com/wiki/index.php?title=Arc}{arc} of a circle is given by the measure of that arc's \href{https://artofproblemsolving.com/wiki/index.php?title=Central\_angle}{central angle} in \href{https://artofproblemsolving.com/wiki/index.php?title=Radians}{radians} times the radius of that circle. Note that the central angle of circle $I$ measures $\frac{\pi}{3}$ radians, and the central angle of circle $II$ measures $\frac{\pi}{4}$ radians. Let the radius of circle $I$ be $r_1$ and the radius of circle $II$ be $r_2$. Then, we know that the arc length $x=\frac{\pi}{3}r_1=\frac{\pi}{4}r_2$, and so $\frac{r_1}{r_2}=\frac{3}{4}$. From this, we see that $\frac{r_1^2}{r_2^2}=\frac{\pi r_1^2}{\pi r_2^2}=\frac{[I]}{[II]}=\frac{9}{16}$. Thus, our answer is $\fbox{(B) 9:16}$.

 


	\item (\textit{AHSME 1968}) A circle passes through the vertices of a triangle with side-lengths $7\frac{1}{2},10,12\frac{1}{2}.$ The radius of the circle is:


 $\text{(A) } \frac{15}{4}\quad \text{(B) } 5\quad \text{(C) } \frac{25}{4}\quad \text{(D) } \frac{35}{4}\quad \text{(E) } \frac{15\sqrt{2}}{2}$

 \textbf{Answer: }C



\textbf{Solution}

The triangle that goes through all the vertices of the triangle is the circumcircle of the triangle.$(7\frac{1}{2})^{2}+10^{2}=(12\frac{1}{2})^{2}$, so the triangle is a right triangle.The radius of a circumcircle of a right triangle is half the hypotenuse. $\frac{1}{2}\cdot \frac{25}{2}=\frac{25}{4}\implies$ $\fbox{C}$



	\item (\textit{AHSME 1968}) If $m$ and $n$ are the roots of $x^2+mx+n=0 ,m \ne 0,n \ne 0$, then the sum of the roots is:


 $\text{(A) } -\frac{1}{2}\quad \text{(B) } -1\quad \text{(C) } \frac{1}{2}\quad \text{(D) } 1\quad \text{(E) } \text{undetermined}$

 \textbf{Answer: }B



\textbf{Solution}

By Vieta's Theorem, $mn = n$ and $-(m + n) = m$. Dividing the first equation by $n$ gives $m = 1$. Multiplying the 2nd by -1 gives $m + n = -m$. The RHS is -1, so the answer is $\fbox{B}$



	\item (\textit{AHSME 1968}) If $x$ and $y$ are non-zero numbers such that $x=1+\frac{1}{y}$ and $y=1+\frac{1}{x}$, then $y$ equals


 $\text{(A) } x-1\quad \text{(B) } 1-x\quad \text{(C) } 1+x\quad \text{(D) } -x\quad \text{(E) } x$

 \textbf{Answer: }E



\textbf{Solution 1}

We see after multiplying the first equation by $y$, that 

 $xy=y+1.$

 Similarly, we see that after multiplying the second equation by $x$, we get that 

 $xy=x+1.$

 Thus $x+1=y+1 \implies x=y$, giving us our final answer of $\fbox{E}.$




\textbf{Solution 2}

We can first combine the fraction on the second equation, which is

 $y=\dfrac{x+1}{x}.$

 Then, substituting the value of $x$ from the first equation, we can get

 $y=\dfrac{1+\dfrac{1}{y}+1}{1+\dfrac{1}{y}}.$

 Then, multiplying the denominator on both sides, we get

 $y+1=2+\dfrac{1}{y} \implies y=1+\dfrac{1}{y}$ which is equivalent to the value of $x$ that was originally provided, giving the final answer of $\fbox{E}.$



	\item (\textit{AHSME 1968}) Let $P$ be the product of any three consecutive positive odd integers. The largest integer dividing all such $P$ is:


 $\text{(A) } 15\quad \text{(B) } 6\quad \text{(C) } 5\quad \text{(D) } 3\quad \text{(E) } 1$

 \textbf{Answer: }D



\textbf{Solution}

Product $P$ can be written as $2n-1$,$2n+1$,$2n+3$. Because $P$ is defined as a "3 consecutive odd integer" product impies that $P$ must be divisible by at least 3. A is ruled out because factors of 5 only arise every 5 terms, if we were to take the 3 terms in the middle of the factors of 5 we wouldn't have a factor of 5. Obviously B is impossible because we are multiplying odd numbers and 2 would never become one of our prime factors. C is ruled out with the same logic as A. Lastly E is ruled out because we have already proved that 3 is possible, and the question asks for greatest possible.$\fbox{D}$

            \begin{itemize}
        \item In first step we could also apply modular arthmetic and get the same answer.
            \end{itemize}
            


	\item (\textit{AHSME 1968}) If $x$ is such that $\frac{1}{x}<2$ and $\frac{1}{x}>-3$, then:


 $\text{(A) } -\frac{1}{3}<x<\frac{1}{2}\quad \text{(B) } -\frac{1}{2}<x<3\quad \text{(C) } x>\frac{1}{2}\quad\\ \text{(D) } x>\frac{1}{2} \text{ or} -\frac{1}{3}<x<0\quad \text{(E) } x>\frac{1}{2} \text{ or } x<-\frac{1}{3}$

 \textbf{Answer: }E



\textbf{Solution}

Because $\frac{1}{x}<2$, $x>\frac{1}{2}$ when $x$ is positive. Because $\frac{1}{x}>-3$, $x<\frac{-1}{3}$ when $x$ is negative. $x$ can never be $0$, so these two inequalities cover all cases for the value of $x$. Thus, our answer is $\fbox{E}$.

 


	\item (\textit{AHSME 1968}) Let $f(n)=\frac{x_1+x_2+\cdots +x_n}{n}$, where $n$ is a positive integer. If $x_k=(-1)^k, k=1,2,\cdots ,n$, the set of possible values of $f(n)$ is:


 $\text{(A) } \{0\}\quad \text{(B) } \{\frac{1}{n}\}\quad \text{(C) } \{0,-\frac{1}{n}\}\quad \text{(D) } \{0,\frac{1}{n}\}\quad \text{(E) } \{1,\frac{1}{n}\}$

 \textbf{Answer: }C



\textbf{Solution}

Because we start with $x_1=-1$, and the terms $x_k$ alternate between $1$ and $-1$, there is either one more $-1$ than the number of $1$s (when $n$ is odd), or there are an equal number of $1$s and $-1$s (when $n$ is even). In the former case, $f(n)=\frac{-1}{n}$, and, in the latter case, $f(n)=0$. This is only consistent with answer $\fbox{C}$.

 


	\item (\textit{AHSME 1968}) Side $AB$ of triangle $ABC$ has length 8 inches. Line $DEF$ is drawn parallel to $AB$ so that $D$ is on segment $AC$, and $E$ is on segment $BC$. Line $AE$ extended bisects angle $FEC$. If $DE$ has length $5$ inches, then the length of $CE$, in inches, is:


 $\text{(A) } \frac{51}{4}\quad \text{(B) } 13\quad \text{(C) } \frac{53}{4}\quad \text{(D) } \frac{40}{3}\quad \text{(E) } \frac{27}{2}$

 \textbf{Answer: }D



\textbf{Solution}

Draw a line passing through $A$ and parallel to $BC$. Let $\angle FEC = 2n$. By alternate-interior-angles or whatever, $\angle BAE = n$, so $BAE$ is an isosceles triangle, and it follows that $BE = 8$. $\triangle ABC \sim \triangle DEC$. Let $CE = x$. We have
 \[\frac{8}{8+x} = \frac{5}{x} \Rightarrow 40+5x = 8x \Rightarrow x = CE =\boxed{\frac{40}{3}}.\]

 Thus, we choose answer $\fbox{D}$.

 


	\item (\textit{AHSME 1968}) Let $n$ be the number of ways $10$ dollars can be changed into dimes and quarters, with at least one of each coin being used. Then $n$ equals:


 $\text{(A) } 40\quad \text{(B) } 38\quad \text{(C) } 21\quad \text{(D) } 20\quad \text{(E) } 19$

 \textbf{Answer: }E



\textbf{Solution}

Because we have to use at least one dime and one quarter, this part of the selection is forced. Thus, the number of ways to partition $\$10$ in the desired manner is equal to the number of ways to partition $\$10-10cent-25cent=\$9.65$ into quarters and dimes without restriction. However, $\$9.65$ is not a multiple of $25cent$, so we must use dimes until we reach a multiple of $25cent$. We first reach such a multiple at $\$9.25$, and, because we have had no freedom up to this point, the number of ways to partition $\$9.25$ into dimes and quarters will yield our desired answer. $\$9.25$ can be partitioned into 37 quarters, and we can replace quarters 2 at a time with 5 dimes to partition it differently. We can replace anywhere from 0 to 18 pairs of quarters with dimes in this manner, giving us $\boxed{19}$ different partitions, which is answer choice $\fbox{E}$.

 


	\item (\textit{AHSME 1968}) The measures of the interior angles of a convex polygon of $n$ sides are in arithmetic progression. If the common difference is $5^{\circ}$ and the largest angle is $160^{\circ}$, then $n$ equals:


 $\text{(A) } 9\quad \text{(B) } 10\quad \text{(C) } 12\quad \text{(D) } 16\quad \text{(E) } 32$

 \textbf{Answer: }A



\textbf{Solution}

The formula for the sum of the angles in any polygon is $180(n-2)$. Because this particular polygon is convex and has its angles in an arithmetic sequence with its largest angle being $160$, we can find the sum of the angles.

 $a_{n}=160$
$a_{1}=160-5(n-1)$

 Plugging this into the formula for finding the sum of an arithmetic sequence...

 $n\left(\frac{160+160-5(n-1)}{2}\right)=180(n-2)$.

 Simplifying, we get $n^2+7n-144$.

 Since we want the positive solution to the quadratic, we can easily factor and find the answer is $n=\boxed{9}$.

 Hence the answer is$\fbox{A}$



	\item (\textit{AHSME 1968}) If $S=1!+2!+3!+\cdots +99!$, then the units' digit in the value of S is:


 $\text{(A) } 9\quad \text{(B) } 8\quad \text{(C) } 5\quad \text{(D) } 3\quad \text{(E) } 0$

 \textbf{Answer: }D



\textbf{Solution}

Note that every factorial after $5!$ has a unit digit of $0$, meaning that we can disregard them. Thus, we only need to find the units digit of $1! + 2! + 3! +  4!$, and as $1! + 2! + 3! +  4! \equiv 3$ mod $10$, which means that the unit digit is $3$, we have our answer of $\fbox{D}$ as desired.

 


	\item (\textit{AHSME 1968}) A segment of length $1$ is divided into four segments. Then there exists a quadrilateral with the four segments as sides if and only if each segment is:


 $\text{(A) equal to } \frac{1}{4}\quad\\ \text{(B) equal to or greater than } \frac{1}{8} \text{ and less than }\frac{1}{2}\quad\\ \text{(C) greater than } \frac{1}{8} \text{ and less than }\frac{1}{2}\quad\\ \text{(D) equal to or greater than } \frac{1}{8} \text{ and less than }\frac{1}{4}\quad\\ \text{(E) less than }\frac{1}{2}$

 \textbf{Answer: }E



\textbf{Solution}

As a consequence of the \href{https://artofproblemsolving.com/wiki/index.php?title=Triangle\_Inequality}{Triangle Inequality}, we can form a quadrilateral from the four segments iff the sum of any three sides is greater than the sum of the fourth side. Thus, no segment can have a length $\geq \frac{1}{2}$. All other segment lengths are valid, so long as the inequality is satsified. Thus, our answer is $\fbox{(E) less than 1/2}$.

 


	\item (\textit{AHSME 1968}) If all the logarithms are real numbers, the equality
$log(x+3)+log(x-1)=log(x^2-2x-3)$
is satisfied for:


 $\text{(A) all real values of }x \quad\\ \text{(B) no real values of } x\quad\\ \text{(C) all real values of } x \text{ except } x=0\quad\\ \text{(D) no real values of } x \text{ except } x=0\quad\\ \text{(E) all real values of } x \text{ except } x=1$

 \textbf{Answer: }B



\textbf{Solution}

From the given we have
 \[\log(x+3)+\log(x-1)=\log(x^2-2x-3)\]
 \[\log(x^2+2x-3)=\log(x^2-2x-3)\]
 \[x^2+2x-3=x^2-2x-3\]
 \[x=0\]However substituing into $\log(x-1)$ gets a negative argument, which is impossible $\boxed{B}$.

 


	\item (\textit{AHSME 1968}) A painting $18$" X $24$" is to be placed into a wooden frame with the longer dimension vertical. The wood at the top and bottom is twice as wide as the wood on the sides. If the frame area equals that of the painting itself, the ratio of the smaller to the larger dimension of the framed painting is:


 $\text{(A) } 1:3\quad \text{(B) } 1:2\quad \text{(C) } 2:3\quad \text{(D) } 3:4\quad \text{(E) } 1:1$

 \textbf{Answer: }C



\textbf{Solution}

Let the width of the frame on the sides to be $x$.

 Then, the width of the frame on the top and bottom is $2x$.

 The area of the frame is then $x\cdot 2x-18\cdot24$

 Setting the area of the frame equal to the area of the picture,
 \[(2x+18)(4x+24)-18\cdot24 = 18\cdot24\]Solving,
 \[8x^2+120x+432 = 864\]
 \[x^2+15x-54=0\]
 \[(x+18)(x-3)=0\Rightarrow x=3\]

 Therefore, the ratio of the smaller side to the larger side is $\frac{18+2\cdot3}{24+4\cdot3}=\frac{24}{36}=\boxed{2:3 (C)}$.

 


	\item (\textit{AHSME 1968}) Ace runs with constant speed and Flash runs $x$ times as fast, $x>1$. Flash gives Ace a head start of $y$ yards, and, at a given signal, they start off in the same direction. Then the number of yards Flash must run to catch Ace is:


 $\text{(A) } xy\quad \text{(B) } \frac{y}{x+y}\quad \text{(C) } \frac{xy}{x-1}\quad \text{(D) } \frac{x+y}{x+1}\quad \text{(E) } \frac{x+y}{x-1}$

 \textbf{Answer: }C



\textbf{Solution}

Let $k$ denotes the distance Ace needs to run after the $y$ yard. Since the distance they run with same amount of time is proportional to their speed, we have
 \[\frac{1}{x}=\frac{k}{y+k}\]
 \[k=\frac{y}{x-1}\]Thus the total distance ran by Flash is
 \[y+k=y+\frac{y}{x-1}=\frac{xy}{x-1}\boxed{C}\]



	\item (\textit{AHSME 1968}) Let $S=2+4+6+\cdots +2N$, where $N$ is the smallest positive integer such that $S>1,000,000$. Then the sum of the digits of $N$ is:


 $\text{(A) } 27\quad \text{(B) } 12\quad \text{(C) } 6\quad \text{(D) } 2\quad \text{(E) } 1$

 \textbf{Answer: }E



\textbf{Solution}

Note that $S = 2(1+2+3+...+N) = N(N +1)$. It follows that $N = 1000$, so the sum of the digits of $N$ is $\fbox{E}$. 

 
FrostFox

 


	\item (\textit{AHSME 1968}) Let $S_n=1-2+3-4+\cdots +(-1)^{n-1}n$, where $n=1,2,\cdots$. Then $S_{17}+S_{33}+S_{50}$ equals:


 $\text{(A) } 0\quad \text{(B) } 1\quad \text{(C) } 2\quad \text{(D) } -1\quad \text{(E) } -2$

 \textbf{Answer: }B



\textbf{Solution}

If $n$ is even, $S_{n}$ is negative $n/2$. If $n$ is odd, then $S_{n}$ is $(n+1)/2$.(These can be found using simple calculations.)Therefore, we know $S_{17}+S_{33}+S_{50}$ =$9+17-25$, which is $\fbox{B}$.

 
Solution By FrostFox

 


	\item (\textit{AHSME 1968}) If the arithmetic mean of $a$ and $b$ is double their geometric mean, with $a>b>0$, then a possible value for the ratio $a/b$, to the nearest integer, is:


 $\text{(A) } 5\quad \text{(B) } 8\quad \text{(C) } 11\quad \text{(D) } 14\quad \text{(E) none of these}$

 \textbf{Answer: }D



\textbf{Solution}

$\fbox{D}$
$\frac{a+b}{2}=2\cdot\sqrt{ab}$
$\frac{a}{b} +1=4\cdot\sqrt{\frac{a}{b}}$

 setting $x=\sqrt{\frac{a}{b}}$ we get a quadratic equation is$x^2+1=4x$ with solutions $x=\frac{4\pm \sqrt{16-4}}{2}$
$x^2=\frac{a}{b}=(4+3)+4\sqrt{3}=13.8=14$.

 


	\item (\textit{AHSME 1968}) Given the three numbers $x,y=x^x,z=x^{x^x}$ with $.9<x<1.0$. Arranged in order of increasing magnitude, they are:


 $\text{(A) } x,z,y\quad \text{(B) } x,y,z\quad \text{(C) } y,x,z\quad \text{(D) } y,z,x\quad \text{(E) } z,x,y$

 \textbf{Answer: }A



\textbf{Solution 1}

Seeing that we need to compare values with exponents, we think \href{https://artofproblemsolving.com/wiki/index.php?title=Logarithms}{logarithms}. Taking the logarithm base $x$ of each term, we obtain $1$, $x$, and $x^x$. Because $0<x<1$, $f(n)=\log_x(n)$ is \href{https://artofproblemsolving.com/wiki/index.php?title=Monotonic}{monotonically decreasing}, so the order of terms by magnitude in our new set of numbers will be reversed compared to the original set (i.e. if $a<b<c$, then $\log_x(a)>\log_x(b)>\log_x(c))$. However, the order of this set will be reversed again (back to the order of the original set) when we take the logarithm base $x$ a second time. After doing this operation, we find the values $0$, $1$, and $x$, which correspond to $x$, $y$, and $z$, respectively. Because $0.9<x<1$, $0<x<1$, and so, by the correspondence detailed above, $x<z<y$, which yields us answer choice $\fbox{A}$.

 


 



\textbf{Solution 2}

Because $0<x<1$, taking $x$ to the $x$th power will bring it closer to $1$, thereby raising its value. Because we have established that $x<x^x$, and $f(n)=x^n$ is a \href{https://artofproblemsolving.com/wiki/index.php?title=Monotonic}{monotonically decreasing} function, we know that $x^x>x^{x^x}$. However, because $x^x<1$, $x^{x^x}$, compared to $x$, will be closer to (but still less than) $1$. Thus, $x^{x^x}>x$. Putting this all together, we see that $x<x^{x^x}<x^x$, or $x<z<y$, which is answer choice $\fbox{A}$.

 


 


	\item (\textit{AHSME 1968}) Convex polygons $P_1$ and $P_2$ are drawn in the same plane with $n_1$ and $n_2$ sides, respectively, $n_1\le n_2$. If $P_1$ and $P_2$ do not have any line segment in common, then the maximum number of intersections of $P_1$ and $P_2$ is:


 $\text{(A) } 2n_1\quad \text{(B) } 2n_2\quad \text{(C) } n_1n_2\quad \text{(D) } n_1+n_2\quad \text{(E) } \text{none of these}$

 \textbf{Answer: }A



\textbf{Solution 1}

Notice how $P_2$ can pass through each line segment of $P_1$ at most twice. To have more than two intersections, the line passing through $P_1$ would have a zigzag shape which is impossible for convex polygons. Therefore, the intersections does not depend on $P_2$ and the answer is $\fbox{A}$




\textbf{Solution 2}

Try to get the answer experimentally. Draw two of the simplest shapes: a square and a triangle and maximize the number of intersections. You will discover it is 6, and the only expression provided that will give 6 is $\fbox{A}$

 Note : this solution isn’t risk free, as the option of “None of These” is given as well.

 


	\item (\textit{AHSME 1968}) 
                    \begin{figure}[H]
                        \centering
                        \includegraphics[width=0.5\textwidth]{Problem/AHSME_I_1968_31_Problem_Figure_1.png}
                    \end{figure}
                    


 In this diagram, not drawn to scale, Figures $I$ and $III$ are equilateral triangular regions with respective areas of $32\sqrt{3}$ and $8\sqrt{3}$ square inches. Figure $II$ is a square region with area $32$ square inches. Let the length of segment $AD$ be decreased by $12\frac{1}{2}$ \% of itself, while the lengths of $AB$ and $CD$ remain unchanged. The percent decrease in the area of the square is:


 $\text{(A)}\ 12\frac{1}{2}\qquad\text{(B)}\ 25\qquad\text{(C)}\ 50\qquad\text{(D)}\ 75\qquad\text{(E)}\ 87\frac{1}{2}$

 \textbf{Answer: }D



\textbf{Solution}

Given an equilateral triangle with side length $s$, the area is given by $\frac{s^2\sqrt{3}}{4}$. Setting this equation equal to the area of triangle $I$, $32\sqrt{3}$, we find that $s=8\sqrt{2}$. Because triangle $III$ is also equilateral, it is similar to trinagle $I$, and, because it has a quarter of the area of $I$, it has $\sqrt{\frac{1}{4}}=\frac{1}{2}$ of the side length. Thus, its sides have a length of $4\sqrt{2}$. Square $II$ initially has an area of $32$, so it's side length starts at $\sqrt{32}=4\sqrt{2}$. The initial length $AD$, therefore, is $16\sqrt{2}$. Because $AD$ decreases by $12.5$ \%, it becomes $\frac{7}{8}$ of its inital value, which is $14\sqrt{2}$. Because the sides of the triangles remain unchanged, this decrease of $2\sqrt{2}$ must come from the side length of the square. Thus, the square's final side is $2\sqrt{2}$, which gives an area of $8$ square inches. $32$ to $8$ is a decrease of $75$ \%. Therefore, our answer is $\boxed{\textbf{(D)}}$.

 


	\item (\textit{AHSME 1968}) $A$ and $B$ move uniformly along two straight paths intersecting at right angles in point $O$. When $A$ is at $O$, $B$ is $500$ yards short of $O$. In two minutes they are equidistant from $O$, and in $8$ minutes more they are again equidistant from $O$. Then the ratio of $A$'s speed to $B$'s speed is:


 $\text{(A) } 4:5\quad \text{(B) } 5:6\quad \text{(C) } 2:3\quad \text{(D) } 5:8\quad \text{(E) } 1:2$

 \textbf{Answer: }C



\textbf{Solution}

Let the speed of $A$ be $a$ and the speed of $B$ be $b$. The first time that $A$ and $B$ will be equidistant from $O$, $B$ will have not yet reached $O$. Thus, after two minutes, $B$'s distance from $O$ will be $500-2b$, and $A$'s distance from $O$ will be $2a$. Setting these expressions equal to each other and dividing by 2, we see that $a=250-b$.

 After another eight minutes (or after a total of ten minutes since $A$ was at $O$), $A$ and $B$ will again be equidistant from $O$, but this time $B$ will have passed $O$. The distance $A$ will be from $O$ is $10a$, and the distance $B$ will be from $O$ is $10b-500$. Setting these expressions equal to each other and dividing by 10, we see that $a=b-50$.

 Adding the two equations that we have obtained above, we see that $2a=250-b+b-50$, and so $a=100$. Substituting this value of $a$ into the second equation, we see that $100=b-50$, or $b=150$. Then, $\frac{a}{b}=\frac{100}{150}=\frac{2}{3}$, so the ratio of $A$'s speed to that of $B$ is $\fbox{(C) 2:3}$.

 


	\item (\textit{AHSME 1968}) A number $N$ has three digits when expressed in base $7$. When $N$ is expressed in base $9$ the digits are reversed. Then the middle digit is:


 $\text{(A) } 0\quad \text{(B) } 1\quad \text{(C) } 3\quad \text{(D) } 4\quad \text{(E) } 5$

 \textbf{Answer: }A



\textbf{Solution}

Call the number $\overline{abc}$ in base 7.

 Then, $49a+7b+c=81c+9b+a$. (Breaking down the number in base-form).

 After combining like terms and moving the variables around,$48a=2b+80c$,$b=40c-24a=8(5c-2a)$. This shows that $b$ is a multiple of 8 (we only have to find the middle digit under \textit{one} of the bases). Thus, $b=0$ (since 8>6, the largest digit in base 7).

 Select $\fbox{A}$ as our answer.

 


	\item (\textit{AHSME 1968}) With $400$ members voting the House of Representatives defeated a bill. A re-vote, with the same members voting, resulted in the passage of the bill by twice the margin by which it was originally defeated. The number voting for the bill on the revote was $\frac{12}{11}$ of the number voting against it originally. How many more members voted for the bill the second time than voted for it the first time?


 $\text{(A) } 75\quad \text{(B) } 60\quad \text{(C) } 50\quad \text{(D) } 45\quad \text{(E) } 20$

 \textbf{Answer: }B



\textbf{Solution}

Let the number of votes for the bill the first time be $a$ and the number of votes for the bill the second time be $b$. Then, from the problem, we know that $b=\frac{12}{11}(400-a)$. Furthermore, the original vote's margin of defeat was $(400-a)-a=400-2a$, and the revote's margin of victory was $b-(400-b)=2b-400$. Thus, from the problem, we know that $2(400-2a)=2b-400$. Substituting the value of $b$ we obtained earlier into this equation, we see that:\begin{align*} \\2(400-2a)&=2*\frac{12}{11}(400-a)-400 \\400-2a&=\frac{400*12}{11}-\frac{12a}{11}-200 \\200-a&=\frac{400*6}{11}-\frac{6a}{11}-100 \\300-\frac{2400}{11}&=\frac{5a}{11} \\5a&=3300-2400 \\5a&=900 \\a&=180 \\\end{align*}
From this fact, we can conclude that $c=\frac{12}{11}*(400-180)=240$, so $c-a=\boxed{60}$, which is answer choice $\fbox{B}$.

 


	\item (\textit{AHSME 1968}) 
                    \begin{figure}[H]
                        \centering
                        \includegraphics[width=0.5\textwidth]{Problem/AHSME_I_1968_35_Problem_Figure_1.png}
                    \end{figure}
                    


 In this diagram the center of the circle is $O$, the radius is $a$ inches, chord $EF$ is parallel to chord $CD$. $O$,$G$,$H$,$J$ are collinear, and $G$ is the midpoint of $CD$. Let $K$ (sq. in.) represent the area of trapezoid $CDFE$ and let $R$ (sq. in.) represent the area of rectangle $ELMF.$ Then, as $CD$ and $EF$ are translated upward so that $OG$ increases toward the value $a$, while $JH$ always equals $HG$, the ratio $K:R$ becomes arbitrarily close to:


 $\text{(A)} 0\quad\text{(B)} 1\quad\text{(C)} \sqrt{2}\quad\text{(D)} \frac{1}{\sqrt{2}}+\frac{1}{2}\quad\text{(E)} \frac{1}{\sqrt{2}}+1$

 \textbf{Answer: }D



\textbf{Solution}

Let $OG = a - 2h$, where $h = JH = HG$. Since the areas of rectangle $EHGL$ and trapezoid $EHGC$ are both half of rectangle $LMFE$ and trapezoid $EFDC$, respectively, the ratios between their areas will remain the same, so let us consider rectangle $EHGL$ and trapezoid $EHGC$.

 Draw radii $OC$ and $OE$, both of which obviously have length $a$. By the Pythagorean Theorem, the length of $EH$ is $\sqrt{a^2 - (OG + h)^2}$, and the length of $CG$ is $\sqrt{a^2 - OG^2}$. It follows that the area of rectangle $EHGL$ is 
 \[EH\cdot HG = h\sqrt{a^2 - (OG + h)^2}\] while the area of trapezoid $EHGC$ is  
 \[\frac{HG}{2}(EH + CG)=\frac{h}{2}\left(\sqrt{a^2 - (OG + h)^2} + \sqrt{a^2 - OG^2}\right).\]

 Now, we want to find the limit, as $OG$ approaches $a$, of $\frac{K}{R}$. Note that this is equivalent to finding the same limit as $h$ approaches $0$. Substituting $a - 2h$ into $OG$ yields that trapezoid $EHGC$ has area 
 \[\frac{h}{2}\left(\sqrt{a^2 - (a - 2h + h)^2} + \sqrt{a^2 - (a - 2h)^2}\right) =\frac{h}{2}\left(\sqrt{2ah - h^2} + \sqrt{4ah - 4h^2}\right)\] and that rectangle $EHGL$ has area 
 \[h\sqrt{a^2 - (a - 2h + h)^2} = h\left(\sqrt{2ah - h^2}\right).\] Our answer thus becomes

 
 \begin{align*} \lim_{h\rightarrow 0}\dfrac{\frac{h}{2}\bigl(\sqrt{2ah - h^2} + \sqrt{4ah - 4h^2}\bigr)}{h\bigl(\sqrt{2ah - h^2}\bigr)} &= \lim_{h\rightarrow 0}\left[\dfrac{1}{2}\cdot\dfrac{\sqrt{h}\bigl(\sqrt{2a - h} + 2\sqrt{a - h}\bigr)}{\sqrt{h}\bigl(\sqrt{2a - h}\bigr)}\right] \\ \implies \frac{1}{2}\cdot\frac{\sqrt{2a} + 2\sqrt{a}}{\sqrt{2a}} &= \frac{1}{2}\left(1 + \frac{2}{\sqrt{2}}\right) = \boxed{\textbf{(D) }\frac{1}{2}+\frac{1}{\sqrt{2}}.} \end{align*}




\end{enumerate}